% Stratégie de protection pour la synthèse peptidique (L'essentiel, chapitre 12)
\begin{tikzpicture}[x=1cm,y=1cm,
  aa/.style={draw=otNavy,fill=white,line width=0.6pt,rounded corners=1.2mm,minimum width=17mm,minimum height=8mm,font=\otMedium\fontsize{8.5}{10}\selectfont,text=otInk},
  cad/.style={draw=otOrange!80!black,fill=otOrange!15,line width=0.6pt,rounded corners=0.8mm,minimum width=7mm,minimum height=6mm,font=\otLight\fontsize{7}{8}\selectfont,text=otOrange!80!black}]
\node[aa] (A) at (0,2.1) {A};
\node[pcLegende,text=otGreenDark,anchor=east] at (-0.95,2.1) {\ce{H2N-}};
\node[pcLegende,text=otNavy,anchor=west] at (0.95,2.1) {\ce{-COOH}};
\node[aa] (B) at (0,0.7) {B};
\node[pcLegende,text=otGreenDark,anchor=east] at (-0.95,0.7) {\ce{H2N-}};
\node[pcLegende,text=otNavy,anchor=west] at (0.95,0.7) {\ce{-COOH}};
\node[cad] at (-1.35,2.75) {P$_1$};
\draw[pcRepere] (-1.35,2.52) -- (-1.1,2.28);
\node[cad] at (1.35,1.35) {P$_2$};
\draw[pcRepere] (1.35,1.12) -- (1.15,0.88);
\node[pcLegendeG,anchor=north,align=center,text width=34mm] at (0.1,-0.1) {1. protéger l'amine de A\\et le carboxyle de B};
\draw[pcFleche] (2.8,1.4) -- (4.1,1.4);
\node[aa,minimum width=34mm] (AB) at (6.05,1.4) {A--CO--NH--B};
\node[pcLegendeG,anchor=north,align=center,text width=36mm] at (6.05,0.75) {2. couplage : seul\\l'enchaînement voulu se forme};
\draw[pcFleche] (8.1,1.4) -- (9.4,1.4);
\node[aa,minimum width=30mm,draw=otGreenDark,text=otGreenDark] (fin) at (11.15,1.4) {dipeptide A--B};
\node[pcLegendeG,anchor=north,align=center,text width=30mm] at (11.15,0.75) {3. déprotéger};
\end{tikzpicture}
\par\smallskip\legende{Les groupes protégés (P$_1$, P$_2$) sont rendus inertes pendant le couplage, puis libérés : un seul dipeptide se forme.}
